\providecommand{\fig}[1]{\figurename~\ref{#1}}

\chapter{Introducción}
\glsresetall
\label{cap:Introduccion}

Este capítulo presenta el contexto general, la problemática, la justificación, el alcance y los objetivos del presente trabajo de titulación, el cual se orienta al diseño e implementación de un mecanismo de cifrado y autenticación transparente para proteger la comunicación entre el controlador de estación \gls{API} y el servidor \gls{SCADA} de la Micro-Red de la Universidad de Cuenca.

\section{Antecedentes}
\label{sec:antecedentes}

La protección de infraestructuras críticas se ha convertido en una prioridad debido a la creciente dependencia tecnológica de sectores como la energía, la industria, la salud y los servicios esenciales. En estos entornos, los \gls{ICS} cumplen un papel central en la supervisión y operación de procesos físicos. Por esta razón, una vulnerabilidad en sus mecanismos de comunicación puede comprometer no solo la información transmitida, sino también la continuidad operativa y la seguridad física de los sistemas controlados \cite{alcaraz2015}.


\section{Identificación del problema}
\label{sec:problema}

La arquitectura de supervisión y control de la Micro-Red de la Universidad de Cuenca depende de la comunicación entre el servidor \gls{SCADA} y el controlador de estación \gls{API}. Este enlace permite el intercambio de información operativa, incluyendo mediciones eléctricas, señales de estado, alarmas y datos asociados al funcionamiento de la infraestructura experimental.


\section{Justificación}
\label{sec:justificacion}

La implementación de un canal seguro transparente se justifica por la necesidad de proteger la información operativa que circula entre el servidor \gls{SCADA} y el controlador de estación \gls{API}. Aunque la red del laboratorio cuenta con mecanismos de segmentación y controles de seguridad perimetral, estos no protegen directamente el contenido de los datos una vez que transitan por la red interna.


\section{Alcance}
\label{sec:alcance}

El presente trabajo se enfoca en el diseño, implementación y validación de un mecanismo de cifrado y autenticación transparente para proteger la comunicación entre el servidor \gls{SCADA} y el controlador de estación \gls{API} de la Micro-Red de la Universidad de Cuenca.


\section{Objetivos}
\label{sec:objetivos}

\subsection{Objetivo general}
Implementar un mecanismo de cifrado y autenticación transparente en el canal de comunicación entre el controlador de estación (\gls{API}) y el servidor \gls{SCADA} de la Micro-Red de la Universidad de Cuenca, para garantizar la confidencialidad e integridad de la telemetría operativa frente a amenazas internas, validando su impacto en el rendimiento y disponibilidad del sistema.

\subsection{Objetivos específicos}
\begin{itemize}
    \item Analizar las tecnologías de cifrado y protocolos de túneles seguros disponibles en el estado del arte, seleccionando la solución que ofrezca el mejor equilibrio entre robustez criptográfica y baja latencia para el protocolo \gls{MODBUSTCP} en el entorno de la Micro-Red.

\end{itemize}